% Pavages de Smith et hypertexte bidimensionnel. I. Arbres
% Compilation : latexmk (voir latexmkrc) — XeLaTeX ou LuaLaTeX.
\documentclass[11pt,a4paper]{article}

\usepackage{fontspec}
\usepackage{polyglossia}
\setdefaultlanguage{french}
\usepackage[margin=2.6cm]{geometry}
\usepackage{amsmath,amssymb,amsthm}
\usepackage{booktabs}
\usepackage{fancyvrb}
\usepackage{caption,subcaption}
\usepackage{enumitem}
\usepackage{csquotes}
\usepackage{microtype}
\usepackage[hidelinks]{hyperref}
\usepackage{smith}

\captionsetup{font=small,labelfont=bf}
\setlength{\emergencystretch}{1.5em}
\newtheorem{theoreme}{Théorème}
\newtheorem{lemme}{Lemme}
\theoremstyle{remark}
\newtheorem*{remarque}{Remarque}

\newcommand{\urlanimation}{https://pgradit.github.io/2DHypertext/papier-I-arbres/deformation.html}
\newcommand{\balise}[1]{\texttt{#1}}
\newcommand{\fig}[1]{\input{figures/#1}}

\title{Pavages de Smith et hypertexte bidimensionnel\\[.3em]
       \large I.~Arbres}
\author{Pierre Gradit}
\date{\today}

\begin{document}
\maketitle

\begin{abstract}
Un document XML bien formé se plonge dans un pavage du plan par neuf tuiles
carrées. Deux lectures de ce pavage, la version graphe et la version disque,
sont duales l'une de l'autre, et une seule règle locale, \papin, appliquée tuile
par tuile, fait passer de l'une à l'autre. La version disque est un disque de
Poincaré à géodésiques intérieures injectives, forme canonique du document ; une
isotopie du plan la relie au pavage.
\end{abstract}

% =========================================================================
\section*{Introduction}

Le projet Xanadu de Ted Nelson concevait l'hypertexte comme une structure de
liens et de transclusions entre documents. Nous en cherchons une version
bidimensionnelle. L'idée directrice est d'exploiter une notion topologique
d'intérieur et d'extérieur, pour obtenir une architecture sûre par construction
(\emph{security-by-design}).

Le point de départ est le pavage de Truchet, tel que C.~S.~Smith l'a
réinterprété avec des quarts de cercle. Nous lui ajoutons une tuile marquée,
\CONNECT, qui porte les arêtes. Les tuiles suffisent à représenter un arbre
comme une carte combinatoire de dimension~2, avec son orientation.

Ce papier se limite aux arbres. Les géodésiques extérieures, le double disque et
les cartes planaires générales sont l'objet du papier~II. Le texte des éléments
n'est pas traité ici : il est porté par la courbe rouge, et sa réécriture au
caractère relève du papier compagnon sur le langage de messagerie.

Le plan suit l'ordre dans lequel un lecteur découvre l'objet. On présente les
tuiles (§\ref{s:tuiles}), puis on plonge l'XML comme un humain le dessinerait
spontanément, sous forme de graphe (§\ref{s:graphe}). On lit ce dessin comme une
carte combinatoire (§\ref{s:carte}), on introduit la dualité (§\ref{s:dualite})
et sa forme locale \papin\ (§\ref{s:papin}), puis on obtient la version disque
(§\ref{s:disque}). La partie se conclut par l'existence d'une déformation
continue entre pavage et disque (§\ref{s:deformation}).

% =========================================================================
\section{Les tuiles}\label{s:tuiles}

Neuf tuiles carrées, assemblées selon une règle de Wang, suffisent à tout le
papier. Elles dérivent de la tuile de Smith.

\subsection{De Truchet à Smith}

La tuile de Truchet est un carré coupé selon une diagonale en deux triangles de
couleurs différentes. C.~S.~Smith remplace cette coupe par deux quarts de cercle.
Chacun relie les milieux de deux côtés adjacents et a pour centre le coin qu'ils
encadrent.

Les deux arcs sont centrés en deux coins opposés. Il y a donc deux orientations,
et un bit par tuile. Chaque arc coupe les côtés à angle droit, en leur milieu.

Une tuile de Smith est ainsi un petit disque à quatre points sur son bord. Ses
deux orientations sont les deux appariements non croisés de ces quatre points ;
le troisième appariement, la croix, est exclu.

\subsection{Les neuf tuiles}

On appelle \emph{courbes rouges} les arcs et segments portés par les tuiles, et
\emph{barre noire} le segment propre à \CONNECT\ (figure~\ref{f:tuiles}).

\begin{figure}[ht]
  \centering
  \fig{fig-tuiles}
  \caption{Les neuf tuiles : \CONNECT\ (2 rotations), \DROIT\ (2), \VIDE\ (1)
    et \VIRAGE\ (4).}
  \label{f:tuiles}
\end{figure}

\begin{center}
\begin{tabular}{@{}llll@{}}
  \toprule
  Tuile & Rotations & Côtés traversés & Contenu \\
  \midrule
  \CONNECT & 2 & 4 & tuile de Smith et barre noire courte \\
  \VIRAGE  & 4 & 2, adjacents & un quart de cercle \\
  \DROIT   & 2 & 2, opposés & un segment entre milieux opposés \\
  \VIDE    & 1 & 0 & rien \\
  \bottomrule
\end{tabular}
\end{center}

La barre noire de \CONNECT\ est portée par la diagonale qui joint les centres des
deux arcs. Elle chevauche les deux arcs un peu au-delà de leurs points de
croisement, puis s'arrête avant les coins : elle reste intérieure à la tuile.

\subsection{Règle d'assemblage}

Chaque côté porte une couleur parmi deux : son milieu est traversé par une
courbe rouge, ou il ne l'est pas. Deux tuiles voisines doivent porter la même
couleur sur leur côté commun. Les neuf tuiles forment ainsi un jeu de tuiles de
Wang.

Cette seule règle garantit que les courbes rouges se prolongent d'une tuile à
l'autre. Dans tout assemblage, elles sont simples, lisses et deux à deux
disjointes. Chacune est fermée, ou se termine au bord de la région pavée.

\subsection{Lecture par les coins}

Marquons chaque coin de tuile « dedans » ou « dehors ». Une courbe rouge sépare
deux coins voisins exactement quand leurs marques diffèrent. Le type de chaque
tuile se lit alors sur ses quatre coins :

\begin{center}
\begin{tabular}{@{}ll@{}}
  \toprule
  Coins « dedans » & Tuile \\
  \midrule
  aucun, ou les quatre & \VIDE \\
  un seul, ou trois & \VIRAGE \\
  deux adjacents & \DROIT \\
  deux opposés & \CONNECT \\
  \bottomrule
\end{tabular}
\end{center}

Le marquage détermine tout le pavage, sauf l'orientation des tuiles \CONNECT. Ce
bit restant dit lequel des deux couples de coins opposés est relié par la bande
comprise entre les arcs. C'est lui qui portera la dualité (§\ref{s:papin}).

% =========================================================================
\section{Plongement de l'XML : la version graphe}\label{s:graphe}

Un arbre XML se plonge dans le pavage en donnant à chaque nœud une boucle rouge
et à chaque lien parent–enfant une tuile \CONNECT. C'est le dessin qu'un humain
fait spontanément, par culture de la théorie des graphes, consciente ou non.

\subsection{L'exemple jouet}

Nous suivons un seul exemple tout au long du papier. Le texte entre balises est
ignoré.

\begin{Verbatim}[xleftmargin=2em,fontsize=\small]
<html>
  <head><title></title></head>
  <body><h1></h1><p></p></body>
</html>
\end{Verbatim}

Son arbre a sept nœuds : le document, racine implicite, puis \balise{html},
\balise{head}, \balise{title}, \balise{body}, \balise{h1} et \balise{p}. Il a six
liens parent–enfant, un par élément.

\subsection{Construction}

On réserve un coin « dedans » par nœud. Deux nœuds liés occupent deux coins
voisins en diagonale. Aucune autre paire de coins « dedans » n'est voisine, ni
par un côté ni en diagonale.

Pour l'exemple, un pavage de $5\times 6$ tuiles suffit. Les coins sont repérés
par (colonne, ligne), de $(0,0)$ en haut à gauche à $(5,6)$.

\begin{center}
\begin{tabular}{@{}lcl@{}}
  \toprule
  Nœud & Coin & Parent \\
  \midrule
  document      & $(1,2)$ & --- \\
  \balise{html}  & $(2,3)$ & document \\
  \balise{head}  & $(3,2)$ & \balise{html} \\
  \balise{title} & $(4,1)$ & \balise{head} \\
  \balise{body}  & $(3,4)$ & \balise{html} \\
  \balise{h1}    & $(4,5)$ & \balise{body} \\
  \balise{p}     & $(2,5)$ & \balise{body} \\
  \bottomrule
\end{tabular}
\end{center}

La lecture par les coins détermine alors chaque tuile. Chaque lien donne une
tuile \CONNECT. On l'oriente de sorte que ses arcs entourent les deux coins
« dedans » : la bande entre les arcs joint les deux coins « dehors », et la barre
noire va d'une boucle à l'autre.

\subsection{Résultat}

Chaque nœud est entouré d'une petite boucle rouge fermée
(figure~\ref{f:graphe}). Chaque barre noire relie la boucle du parent à celle de
l'enfant, et porte le nom de l'enfant : \balise{html} relie le document à
\balise{html}, \balise{head} relie \balise{html} à \balise{head}, et ainsi de
suite.

\begin{figure}[ht]
  \centering
  \fig{fig-graphe}
  \caption{L'exemple jouet en version graphe : sept boucles rouges, six tuiles
    \CONNECT. Le spin du vide \spin\ occupe la tuile \VIDE\ en haut à gauche.}
  \label{f:graphe}
\end{figure}

Le pavage compte 30 tuiles : 6 \CONNECT, 16 \VIRAGE, 8 \VIDE\ et aucune \DROIT.
La tuile \DROIT\ apparaît dès qu'un nœud occupe plusieurs coins, par exemple pour
loger son texte.

Les noms sont posés parallèlement aux barres, à $+45^\circ$ ou $-45^\circ$, dans
le sens qui se lit, à écart constant de la barre. Le texte des éléments, lui,
est porté par les courbes rouges ; il relève du papier compagnon.

% =========================================================================
\section{Carte combinatoire et spin du vide}\label{s:carte}

Le pavage en version graphe se lit comme une carte combinatoire orientée de
dimension~2. L'orientation est fixée par un seul symbole, \spin\ (U+2941), posé
dans une tuile \VIDE\ extérieure.

\subsection{Rappels}

Une carte combinatoire est un triplet $(D,\sigma,\alpha)$. $D$ est un ensemble
fini de brins, ou demi-arêtes. $\sigma$ est une permutation de $D$ dont les
cycles sont les sommets. $\alpha$ est une involution sans point fixe dont les
cycles sont les arêtes. Les faces sont les cycles de la permutation dérivée
\[
  \varphi = \sigma\circ\alpha .
\]
En notant $V$, $E$ et $F$ les nombres de cycles de $\sigma$, $\alpha$ et
$\varphi$, une carte connexe vérifie la formule d'Euler
\[
  V - E + F = 2 - 2g ,
\]
et elle est planaire quand le genre $g$ est nul. La notation de Lienhardt,
reprise par CGAL, écrit la même structure avec $\beta_1=\varphi$ et
$\beta_2=\alpha$.

\subsection{Le spin du vide}

Une carte combinatoire est toujours orientée : $\sigma$ tourne dans un sens, et
le sens opposé donne $\sigma^{-1}$, la carte miroir. Le pavage doit donc porter
ce choix. Nous le plaçons dans une tuile \VIDE\ extérieure, sous la forme du
symbole \spin, flèche circulaire fermée dans le sens horaire.

Ce placement a un avantage décisif. La règle \papin\ du §\ref{s:papin} ne
touche que les tuiles \CONNECT : l'orientation reste donc la même avant et après
la dualité.

\subsection{Lecture du pavage}

\begin{center}
\begin{tabular}{@{}ll@{}}
  \toprule
  Carte & Pavage \\
  \midrule
  sommet & boucle rouge \\
  brin & extrémité d'une barre noire, là où elle chevauche une boucle \\
  arête ($\alpha$) & les deux extrémités d'une même barre \\
  rotation ($\sigma$) & ordre des brins le long d'une boucle, dans le sens de \spin \\
  \bottomrule
\end{tabular}
\end{center}

\subsection{L'exemple}

Les douze brins sont numérotés de 0 à 11. La numérotation est choisie pour
coïncider avec l'ordre des balises du document, ce que le §\ref{s:disque}
justifie.
\begin{align*}
  \alpha   &= (0\;11)(1\;4)(2\;3)(5\;10)(6\;7)(8\;9),\\
  \sigma_g &= (0)(1\;5\;11)(2\;4)(3)(6\;8\;10)(7)(9).
\end{align*}
Chaque arête a un brin sur la boucle du parent et un sur celle de l'enfant. Par
exemple, la barre \balise{head} a le brin~1 sur la boucle de \balise{html} et le
brin~4 sur celle de \balise{head}. Autour de \balise{html}, dans le sens horaire,
on rencontre \balise{head} (brin~1), \balise{body} (brin~5) puis le document
(brin~11).

Le calcul de $\varphi$ donne une seule face, le tour complet de l'arbre :
\[
  \varphi_g = \sigma_g\circ\alpha = (0\;1\;2\;3\;4\;5\;6\;7\;8\;9\;10\;11).
\]
On vérifie $7-6+1=2$ : la carte est planaire, et c'est un arbre plan.

% =========================================================================
\section{Dualité}\label{s:dualite}

La carte duale échange sommets et faces en gardant les mêmes brins et les mêmes
arêtes. Appliquée à la version graphe, elle donne une carte à un seul sommet et
sept faces.

\subsection{Définition}

La duale de $(D,\sigma,\alpha)$ est $(D,\varphi,\alpha)$. Avec la convention
$\varphi=\sigma\circ\alpha$, la face de la duale est
$\varphi\circ\alpha=\sigma$, puisque $\alpha$ est une involution. Passer deux fois
au dual ramène donc exactement à la carte de départ :
\[
  (D,\ \sigma,\ \alpha)\ \longleftrightarrow\ (D,\ \sigma\circ\alpha,\ \alpha).
\]
L'orientation est conservée : $\sigma$ et $\varphi$ tournent dans le sens fixé
par \spin.

\subsection{L'exemple}

La duale de la version graphe a pour rotation $\sigma_d=\varphi_g$, un seul cycle
de douze brins. Ses faces sont les cycles de $\varphi_d=\sigma_d\circ\alpha$, qui
redonnent les sept boucles de la version graphe.

\begin{center}
\begin{tabular}{@{}lll@{}}
  \toprule
  & Version graphe & Version duale \\
  \midrule
  rotation $\sigma$ & $(0)(1\,5\,11)(2\,4)(3)(6\,8\,10)(7)(9)$ & $(0\,1\,2\,\dots\,11)$ \\
  arêtes $\alpha$ & $(0\,11)(1\,4)(2\,3)(5\,10)(6\,7)(8\,9)$ & les mêmes \\
  faces $\varphi$ & $(0\,1\,2\,\dots\,11)$ & $(0)(1\,5\,11)(2\,4)(3)(6\,8\,10)(7)(9)$ \\
  $V-E+F$ & $7-6+1=2$ & $1-6+7=2$ \\
  \bottomrule
\end{tabular}
\end{center}

Les sept faces de la duale sont les régions délimitées par les arêtes autour de
l'unique sommet. Elles correspondent une à une aux nœuds de l'arbre : $(0)$ au
document, $(1\,5\,11)$ à \balise{html}, $(2\,4)$ à \balise{head}, $(3)$ à
\balise{title}, $(6\,8\,10)$ à \balise{body}, $(7)$ à \balise{h1} et $(9)$ à
\balise{p}.

Reste à réaliser cette duale dans le pavage, sans changer de jeu de tuiles.
C'est l'objet de la règle \papin.

% =========================================================================
\section{\texorpdfstring{\papin}{(pa)(pin)}, atome de dualité}\label{s:papin}

La règle \papin\ change l'orientation d'une seule tuile \CONNECT. Elle agit sur
la carte comme une transposition, et l'appliquer à toutes les tuiles \CONNECT\
en un seul pas réalise le passage au dual.

\subsection{La règle}

Dans une tuile \CONNECT, les arcs changent de coins et la barre noire change de
diagonale (figure~\ref{f:papin}). La bande qui joignait deux coins opposés joint
désormais les deux autres. Le nom de l'élément reste attaché à la barre.

\begin{figure}[ht]
  \centering
  \fig{fig-papin}
  \caption{La règle \papin\ sur une tuile \CONNECT. La lettre \emph{p} est le
    nom générique de la barre noire.}
  \label{f:papin}
\end{figure}

Les quatre côtés restent traversés dans les deux orientations. La règle de Wang
est donc préservée : \papin\ est une règle de réécriture locale, qui transforme
un pavage valide en un pavage valide.

Le nom vient du robinet de Denis Papin, qui fonctionne exactement sur ce
principe. La récurrence du \emph{p} fait de \emph{p} le nom générique posé sur
l'arête noire de \CONNECT.

\subsection{Effet sur la carte}

Soit une barre d'extrémités $x$ et $y$, donc une arête $(x\;y)$ de $\alpha$.
Appliquer \papin\ à sa tuile laisse $\alpha$ inchangée et remplace la rotation
par
\[
  \sigma' = \sigma\circ(x\;y).
\]
Composer par une transposition change le nombre de cycles de $\pm1$. Dans le
pavage, la règle soit fusionne deux boucles rouges en une, soit en coupe une en
deux. Appliquée deux fois, elle ramène la tuile à son état initial.

\subsection{Toutes les tuiles à la fois}

Les arêtes de $\alpha$ sont des transpositions disjointes : elles commutent et
leur produit est $\alpha$. Appliquer \papin\ à toutes les tuiles \CONNECT, dans
n'importe quel ordre, donne
\[
  \sigma\circ\prod_{(x\,y)\in\alpha}(x\;y) = \sigma\circ\alpha = \varphi .
\]
La nouvelle rotation est l'ancienne face : c'est exactement la dualité du
§\ref{s:dualite}. Le symbole \spin\ n'est pas touché, donc l'orientation est la
même des deux côtés.

\subsection{L'exemple}

Partant de la version graphe, chacun des six retournements fusionne deux
boucles. On passe de 7 boucles rouges à 1, et d'une face à 7
(figure~\ref{f:versions}). Les sept petites boucles des nœuds ont disparu : il
reste une boucle unique, et sept régions séparées par les barres.

\begin{figure}[ht]
  \centering
  \begin{subfigure}{.45\textwidth}\centering
    \fig{fig-graphe}\caption{Version graphe}
  \end{subfigure}\hfill
  \begin{subfigure}{.45\textwidth}\centering
    \fig{fig-disque-pavage}\caption{Version disque}
  \end{subfigure}
  \caption{\papin\ appliquée aux six tuiles \CONNECT\ : de sept boucles à une.}
  \label{f:versions}
\end{figure}

% =========================================================================
\section{Version disque : Poincaré et géodésiques injectives}\label{s:disque}

Après \papin\ sur toutes les tuiles \CONNECT, l'unique boucle rouge est l'horizon
d'un disque de Poincaré. Les barres noires en sont les géodésiques intérieures,
sans autre intersection que l'horizon. C'est la forme canonique d'un XML.

\subsection{Le disque}

On place $2n$ points ordonnés sur le cercle unité, l'horizon. Chaque point est
l'extrémité d'une et une seule arête. Chaque arête est un arc de cercle
orthogonal à l'horizon, ou un diamètre si ses extrémités sont antipodales.

Les arêtes sont \emph{injectives} : elles ne se croisent pas, et ne rencontrent
l'horizon qu'en leurs extrémités. Pour figurer le chevauchement des barres de
\CONNECT, chaque géodésique déborde légèrement de l'horizon, perpendiculairement
au cercle (figure~\ref{f:poincare}).

\begin{figure}[ht]
  \centering
  \fig{fig-poincare}
  \caption{L'exemple jouet en disque de Poincaré : six géodésiques intérieures
    injectives. Les positions $0$ à $11$ se lisent dans le sens de \spin\ à
    partir du haut.}
  \label{f:poincare}
\end{figure}

\subsection{Le lien canonique avec XML}

Les douze brins de l'exemple sont les douze points de l'horizon, parcourus dans
le sens de \spin\ à partir du haut. La rotation $\sigma_d=(0\,1\,2\,\dots\,11)$
est l'ordre de lecture des balises. L'involution $\alpha$ apparie chaque balise
ouvrante à sa fermante.

\begin{center}
\begin{tabular}{@{}lc@{}}
  \toprule
  Balises & Positions \\
  \midrule
  \balise{<html>} \dots\ \balise{</html>}   & $0$ -- $11$ \\
  \balise{<head>} \dots\ \balise{</head>}   & $1$ -- $4$ \\
  \balise{<title>} \dots\ \balise{</title>} & $2$ -- $3$ \\
  \balise{<body>} \dots\ \balise{</body>}   & $5$ -- $10$ \\
  \balise{<h1>} \dots\ \balise{</h1>}       & $6$ -- $7$ \\
  \balise{<p>} \dots\ \balise{</p>}         & $8$ -- $9$ \\
  \bottomrule
\end{tabular}
\end{center}

Pour une carte à un seul sommet, être planaire équivaut à avoir des arêtes qui ne
se croisent pas. Or un appariement non croisé de positions est exactement un mot
bien parenthésé, c'est-à-dire une suite de balises correctement imbriquées. Un
XML bien formé à $n$ éléments est donc un disque à $2n$ points et $n$ géodésiques
intérieures injectives, et réciproquement.

Ces objets sont comptés par le nombre de Catalan $C_n$, comme les arbres plans à
$n$ arêtes. Pour l'exemple, $C_6=132$.

\subsection{Les régions}

Les six géodésiques découpent le disque en sept régions, les faces $\varphi_d$
du §\ref{s:dualite}. La région bordée par l'arc de l'horizon entre 11 et 0 est
celle du document.

Chaque nom est posé du côté de la région que son élément délimite. Pour
\balise{html}, ce côté est vers l'intérieur du disque, puisque le côté horizon
appartient au document.

% =========================================================================
\section{Existence d'une déformation continue}\label{s:deformation}

Le pavage en version disque et le disque de Poincaré du §\ref{s:disque} sont
reliés par une isotopie du plan. La preuve combine le théorème de
Jordan--Schoenflies et un lemme de redressement des cordes, démontré par
récurrence. Une animation de cette correspondance accompagne le
papier\,\footnote{\url{\urlanimation}}.

\subsection{Énoncé}

Notons $\Gamma$ l'unique boucle rouge du pavage, et $b_1,\dots,b_n$ les barres
noires restreintes à leur portion comprise entre leurs deux points de
chevauchement avec $\Gamma$.

\begin{theoreme}\label{t:deformation}
Il existe une isotopie $(h_t)_{t\in[0,1]}$ d'homéomorphismes du plan, avec
$h_0=\mathrm{id}$, telle que $h_1(\Gamma)$ soit l'horizon et que chaque
$h_1(b_k)$ soit la géodésique de même paire. L'orientation fixée par \spin\ est
préservée.
\end{theoreme}

\subsection{Préliminaires}

\begin{enumerate}[label=(\roman*)]
  \item \emph{$\Gamma$ est une courbe de Jordan.} C'est une boucle unique
    (§\ref{s:papin}), simple car les courbes rouges sont disjointes
    (§\ref{s:tuiles}), et lisse par morceaux.
  \item \emph{Les barres sont des cordes intérieures.} Une courbe rouge sépare
    deux coins exactement quand leurs marques diffèrent, et le bord de la région
    pavée est « dehors ». Les coins « dedans » sont donc dans le domaine borné
    $U$ de $\Gamma$. Après \papin, la bande de chaque \CONNECT\ joint deux coins
    « dedans », donc chaque $b_k$ est un arc simple intérieur à $U$, d'extrémités
    sur $\Gamma$.
  \item \emph{Les barres sont disjointes}, chacune étant intérieure à sa tuile.
  \item \emph{L'ordre des extrémités} le long de $\Gamma$, dans le sens de
    \spin, est $\sigma_d$, et leur appariement est $\alpha$.
\end{enumerate}

\subsection{Lemme de redressement}

\begin{lemme}\label{l:redressement}
Soient deux familles de $n$ cordes disjointes d'un disque fermé, ayant les mêmes
extrémités et le même appariement. Il existe un homéomorphisme du disque, fixant
son bord point par point, qui envoie la première famille sur la seconde.
\end{lemme}

\begin{proof}
Par récurrence sur $n$. Pour $n=0$, l'identité convient. Pour $n\ge1$, un
appariement non croisé contient une paire d'extrémités consécutives sur le bord ;
soit $c$ la corde correspondante dans la première famille, et $c'$ dans la
seconde.

Chacune découpe le disque en deux disques fermés, et le petit côté ne contient
aucune autre corde. On fixe un homéomorphisme de $c$ sur $c'$ qui fixe les
extrémités. Avec l'identité sur le bord, il définit un homéomorphisme entre les
bords de chaque côté, que Schoenflies étend aux disques. On obtient $H_1$, qui
fixe le bord et envoie $c$ sur $c'$.

Les $n-1$ autres cordes, transportées par $H_1$, et celles de la seconde famille
sont alors des cordes du grand côté de $c'$, avec les mêmes extrémités et le même
appariement. L'hypothèse de récurrence donne $H_2$ sur ce côté, fixant son bord,
que l'on prolonge par l'identité sur le petit côté. $H_2\circ H_1$ convient.
\end{proof}

\subsection{Preuve du théorème}

\begin{proof}[Preuve du théorème~\ref{t:deformation}]
\leavevmode
\begin{enumerate}[label=\arabic*.]
  \item \emph{Arrondir $\Gamma$.} Par Schoenflies, un homéomorphisme $F$ du plan
    envoie $\Gamma$ sur le cercle unité et $U$ sur le disque ouvert. Quitte à le
    composer avec une symétrie, $F$ préserve l'orientation, donc le sens de
    \spin.
  \item \emph{Répartir les extrémités.} $F$ envoie les $2n$ extrémités sur le
    cercle, dans le même ordre cyclique. Un homéomorphisme croissant du cercle
    les envoie sur les $2n$ points réguliers de l'horizon ; on le prolonge
    radialement en un homéomorphisme $G$ du plan.
  \item \emph{Redresser les cordes.} Les images $G\circ F(b_k)$ et les
    géodésiques forment deux familles de cordes disjointes, de mêmes extrémités
    et de même appariement $\alpha$. Le lemme~\ref{l:redressement} donne un
    homéomorphisme $H$ du disque fixant le cercle, prolongé par l'identité à
    l'extérieur.
  \item \emph{Déformer.} $h=H\circ G\circ F$ est un homéomorphisme du plan qui
    préserve l'orientation et réalise l'énoncé. Tout tel homéomorphisme est
    isotope à l'identité (résultat classique ; référence à fixer), ce qui
    fournit $(h_t)$.\qedhere
\end{enumerate}
\end{proof}

\subsection{Conséquences}

La preuve n'utilise du pavage que $\sigma_d$ et $\alpha$. À isotopie près, le
dessin est donc entièrement déterminé par la suite des balises du document.

\begin{remarque}
L'animation interpole les points en ligne droite entre le pavage et le disque.
Elle montre la correspondance, mais ne garantit pas l'absence de croisement à
chaque instant : c'est l'isotopie $(h_t)$ qui en tient lieu.
\end{remarque}

% =========================================================================
\section{Conclusion et perspectives}

Neuf tuiles de Wang, une orientation posée dans le vide et une règle locale
suffisent à représenter un document XML et sa forme duale. La version graphe est
le dessin spontané de l'arbre. La version disque est sa forme canonique, et une
isotopie du plan relie le pavage au disque de Poincaré.

Le pavage permet aussi des tuiles \CONNECT\ dont la bande est extérieure à
l'horizon. Elles donnent des géodésiques extérieures, donc le double disque, les
mots à deux types de parenthèses et les cartes planaires générales. C'est l'objet
du papier~II, d'où la numérotation de ce titre, à la manière de de~Bruijn pour
ses deux articles sur les pavages de Penrose.

Le texte des éléments, porté par les courbes rouges, est traité dans le papier
compagnon sur le langage de messagerie.

\section*{Références (à vérifier)}
\begingroup\sloppy

\begin{itemize}[leftmargin=*]
  \item S.~Truchet, \emph{Mémoire sur les combinaisons}, Mémoires de l'Académie
    royale des sciences, 1704.
  \item C.~S.~Smith, \emph{The Tiling Patterns of Sebastien Truchet and the
    Topology of Structural Hierarchy}, Leonardo, 1987.
  \item H.~Wang, \emph{Proving theorems by pattern recognition~II}, Bell System
    Technical Journal, 1961.
  \item N.~G.~de~Bruijn, \emph{Algebraic theory of Penrose's non-periodic
    tilings of the plane, I} et \emph{II}, Indagationes Mathematicae, 1981.
  \item P.~Lienhardt, travaux sur les cartes combinatoires et généralisées de
    dimension~$n$.
  \item T.~H.~Nelson, travaux sur le projet Xanadu.
  \item Théorème de Jordan--Schoenflies, et isotopie à l'identité des
    homéomorphismes du plan préservant l'orientation : références à fixer.
\end{itemize}
\endgroup

\end{document}
